\begin{frame}{window=1 vs 2를 말뭉치로 재보기}
  ``사자는''(동물)과 ``바나나는''(과일)의 코사인 유사도
  \textbf{차이} (= ``동물-과일 구분력'')를 window별로 측정:
  \[
    \text{window}=1:\; \text{gap} = +0.303
    \qquad
    \text{window}=2:\; \text{gap} = +0.341
  \]
  \begin{itemize}
    \item window=2가 \textbf{약간} 더 좋은 구분력 --- 이 작은 말뭉치
     에서는 window=2가 ``고양이는 동물이다''처럼 \textbf{세 단어
      패턴}을 포착할 수 있어서(1이면 ``고양이는-동물이다''만, 2면
      ``고양이는-동물이다-강아지는''도) 문맥 정보가 조금 더 담김.
    \item window를 5, 10으로 키워도 \textbf{더 좋은 결과는 안
      나옴} --- 이 작은 말뭉치(고유어 10개, 16 토큰)에선.
    \item 이 실험이 보여주는 것: ``window 최적값''은
      \textbf{말뭉치의 길이와 어휘 크기}에 따라 달라진다.
  \end{itemize}
\end{frame}
